\documentclass[11pt,a4paper]{article}
\usepackage[T1]{fontenc}
\usepackage{lmodern}
\usepackage[margin=27mm]{geometry}
\usepackage{amsmath,amssymb,amsthm,microtype,booktabs}
\usepackage[hidelinks]{hyperref}
\newtheorem{lemma}{Lemma}
\newtheorem{proposition}[lemma]{Proposition}
\theoremstyle{definition}
\newtheorem{definition}[lemma]{Definition}
\newtheorem{example}[lemma]{Example}
\newcommand{\E}{\mathbb E}
\emergencystretch=2em
\clubpenalty=10000
\widowpenalty=10000
\hypersetup{pdftitle={The Soft-Count Measure for Adaptive Branching}}
\title{The Soft-Count Measure for Adaptive Branching\\
\large Definition, one branching step, and examples}
\author{}
\date{8 October 2026}
\begin{document}
\maketitle

\begin{abstract}
We define a complexity measure for the adaptive branching algorithm with
learned pruning constraints and final enumeration. It is a soft count of the
completions that respect the learned constraints. It is exactly additive under
branching, it bounds the remaining work, and a new constraint changes it by an
explicit factor. A constraint has a weight that grows as the constraint
tightens and becomes infinite when the constraint assigns its variables. A new
constraint lowers the measure exactly when, under the current soft count, its
expected number of bad values exceeds its allowance. This is a precise form of
a sufficiently new constraint. We work out one branching step and three small
examples, and run the search on small formulas.
\end{abstract}

\section{The search and its constraints}

The algorithm branches on one free variable at a time, chosen adaptively. When
it branches on $x$ with preferred literal $\ell$, the child in which $\ell$ is
false learns a constraint on the free neighbours of $x$. For each clause
$(\ell\lor m)$ the value of the neighbour that makes $m$ false is \emph{bad},
and for each clause $(\neg\ell\lor m)$ the value that makes $m$ true is bad.
The constraint says that at most $k$ neighbours take their bad value, where
$k=p-1$ and $p$ is the deficit of $\ell$. Indeed, if $p$ or more are bad,
flipping $x$ back to $\ell$ is no worse. Simplification rules such as
domination and propagation may be applied. At the end, the completions
respecting the learned constraints are enumerated.

At a node, a learned constraint $i$ has a \emph{support} $N_i$ (its free
variables), a bad value $\beta_i(w)$ for each $w\in N_i$, and a current
\emph{allowance} $k_i$. For a completion $y$ of the free variables, $L_i(y)$ is
the number of $w\in N_i$ with $y_w=\beta_i(w)$. When a support variable is
assigned it leaves $N_i$, and if it took its bad value, $k_i$ drops by one.
The constraint is violated when $k_i<0$, and \emph{vacuous} when
$k_i\ge|N_i|$.

\section{The measure}

\begin{definition}[Soft count]
Give each constraint a weight $\lambda_i\ge0$. For a free variable $w$ and a
value $b$, let $A_w^b$ be the sum of $\lambda_i$ over the constraints in which
$b$ is the bad value of $w$. The soft count of the node is
\[
 \Phi=\sum_y\ \prod_i e^{\lambda_i(k_i-L_i(y))}
     =e^{\sum_i\lambda_ik_i}\prod_{w\text{ free}}
      \bigl(e^{-A_w^0}+e^{-A_w^1}\bigr),
\]
where $y$ ranges over the completions of the free variables.
\end{definition}

The two expressions agree because $\sum_i\lambda_iL_i(y)=\sum_wA_w^{y_w}$.
Each completion is discounted exponentially in how much it violates each
constraint. Under the \emph{soft distribution}, the distribution on completions
proportional to these terms, the free variables are independent, and $w$
takes the value $b$ with probability
\[
 \pi_w(b)=\frac{e^{-A_w^b}}{e^{-A_w^0}+e^{-A_w^1}} .
\]

\begin{proposition}[Basic properties]\label{prop:basic}
The soft count has three properties.
\begin{enumerate}
 \item $\Phi$ is at least the number of completions that respect every learned
 constraint. If $\Phi<1$, no completion respects them all, and the node can
 be cut.
 \item At the root, $\Phi=2^n$.
 \item \emph{Additivity.} If $x$ is free and the weights are kept, the
 children $x=0$ and $x=1$ have soft counts $\Phi\,\pi_x(0)$ and
 $\Phi\,\pi_x(1)$; together they make up exactly $\Phi$.
\end{enumerate}
\end{proposition}

\begin{proof}
A completion respecting all constraints has $k_i-L_i(y)\ge0$ for every $i$,
so its term is at least $1$. At the root there are no constraints. Fixing
$x=b$ keeps the completions with $y_x=b$; in the product this replaces the
factor of $x$ by $e^{-A_x^b}$, and the constraints in which $b$ is bad lose
one unit of allowance. That gives $\Phi\,\pi_x(b)$.
\end{proof}

\section{One branching step}

Branch on $x$, with preferred value $v$ and other value $\bar v$.
\begin{itemize}
 \item \emph{Preferred child} $x=v$. Every constraint containing $x$ loses
 $x$ from its support, and its allowance drops by one if $v$ is its bad
 value. No constraint is learned. Its soft count is $\Phi\,\pi_x(v)$.
 \item \emph{Other child} $x=\bar v$. The same updates, plus the new
 constraint $c$ on the free neighbours of $x$, with allowance $k_c=p-1$.
 Before $c$ is added, its soft count is $\Phi\,\pi_x(\bar v)$.
\end{itemize}

\begin{lemma}[The factor of a new constraint]\label{lem:factor}
Adding $c$ with weight $\lambda$ multiplies the soft count of the child by
\[
 F_c(\lambda)=e^{\lambda k_c}\prod_{w\in N_c}
 \Bigl(1-q_w\bigl(1-e^{-\lambda}\bigr)\Bigr),
 \qquad q_w=\pi_w(\beta_c(w)),
\]
where $q_w$ is the probability, under the soft distribution of the child,
that $w$ takes its bad value for $c$. In particular $F_c(0)=1$, and the best
weight gives $F_c\le1$.
\end{lemma}

\begin{proof}
The new constraint adds $\lambda$ to $A_w^{\beta_c(w)}$ for $w\in N_c$ and
$\lambda k_c$ to the exponent, so the factor of $w$ becomes
$\pi_w(\text{good})+\pi_w(\text{bad})e^{-\lambda}$.
\end{proof}

The derivative of $\ln F_c$ at $\lambda=0$ is $k_c-\sum_{w\in N_c}q_w$. So:

\begin{quote}
\emph{A new constraint lowers the measure exactly when, under the current soft
count, its expected number of bad values exceeds its allowance.}
\end{quote}

This is the precise version of a sufficiently new constraint. On a fresh
variable $q_w=\frac12$. On a variable where the earlier constraints have the
\emph{same} bad value, $q_w$ is small, so the new constraint adds little: it is
redundant there. Where they have the \emph{opposite} bad value, $q_w$ is
large, so it adds a lot.

The \emph{information} of the step is the decrease of $\log_2\Phi$:
\[
 I_x=-\log_2\Bigl(\pi_x(v)+\pi_x(\bar v)\,F_c\Bigr)\ \ge\ 0 .
\]
The adaptive rule branches on the free variable with the largest $I_x$. A
constraint that becomes vacuous gets weight $0$: its factor was at least $1$,
so this lowers $\Phi$.

\section{The weights}

\paragraph{A single constraint.} With $r$ free variables and allowance
$k<r/2$, on fresh variables, the best weight is
\[
 \lambda=\ln\frac{r-k}{k},
\]
and then the constraint's factor is $2^{rH(k/r)}/2^r$, where $H$ is the binary
entropy. So the weight grows as the constraint tightens. For $k=0$ it is
infinite: the bad values get weight $0$, the constraint assigns all its
variables, and propagation is the limiting case.

\paragraph{Re-optimizing at every node.} We may choose all weights afresh at
every node to minimize $\Phi$. Write $\Phi^*$ for this minimum. For the two
children,
\[
 \Phi^*(x=0)+\Phi^*(x=1)\ \le\ \Phi^*(\text{parent}).
\]
Evaluate both children at the parent's optimal weights, with weight $0$ for
the new constraint. By additivity their soft counts add up to
$\Phi^*(\text{parent})$, and each child's minimum is at most this value. So
re-optimized weights still give a valid measure, and constraints near
assigning their variables automatically receive large weights.
$\ln\Phi$ is convex in the weights, so the minimum can be computed.

\section{Examples}

\begin{example}[One constraint, allowance zero]\label{ex:A}
Let $x_1$ have the clauses $(x_1\lor x_2)$, $(x_1\lor x_3)$,
$(x_1\lor\neg x_4)$ and $(\neg x_1\lor x_5)$. Its preferred value is $1$, and
$p=1$. The child $x_1=0$ learns
\[
 [x_2=0]+[x_3=0]+[x_4=1]+[x_5=1]\le0 .
\]
Here $\Phi=(1+e^{-\lambda})^4$, which is $16$, $3.50$, $1.66$ and $1$ for
$\lambda=0,1,2,\infty$. Exactly one completion is permitted. Branching on
$x_2$ at $\lambda=1$ splits $3.501$ into $2.559$ for $x_2=1$ and $0.942$ for
$x_2=0$. At $\lambda=\infty$ the child $x_2=0$ has soft count $0$ and is cut:
this is propagation.
\end{example}

\begin{example}[Redundant and new constraints]\label{ex:B}
The constraint ``at most one of $y_1,\dots,y_4$ is $1$'' (bad value $1$
everywhere, allowance $1$) is present with weight $\ln3$. Then $\Phi=9.48$,
while $5$ completions are permitted, and $q_w=\frac14$ for bad value $1$. A
second constraint with allowance $1$ arrives on the same variables.

\medskip
\centering
\begin{tabular}{lccccc}
\toprule
bad values of the new constraint & $q_w$ & best $F_c$ & $\Phi$ &
$\Phi^*$ & permitted\\
\midrule
$1,1,1,1$ (agrees) & $\frac14$ each & $1$ & $9.48$ & $9.48$ & $5$\\
$1,1,0,0$ (half) & $\frac14,\frac14,\frac34,\frac34$ & $0.51$ & $4.79$ &
$4.00$ & $2$\\
$0,0,0,0$ (disagrees) & $\frac34$ each & $0.11$ & $1.05$ & $0$ & $0$\\
\bottomrule
\end{tabular}
\par\medskip
\raggedright
The agreeing constraint is redundant and gains nothing. The disagreeing one
makes the node infeasible, and the re-optimized measure $\Phi^*$ detects this.
\end{example}

\begin{example}[A full branching step]\label{ex:C}
Take the formula with the nine clauses
\begin{gather*}
 (x_1\lor x_2),\ (x_1\lor\neg x_3),\ (\neg x_1\lor x_4),\
 (\neg x_2\lor x_3),\ (x_2\lor x_5),\\
 (x_3\lor\neg x_6),\ (x_4\lor\neg x_5),\ (\neg x_4\lor x_6),\
 (x_5\lor x_6).
\end{gather*}
At the root, $x_1$ has gap $+1$ over three clauses, so its preferred value is
$1$ and $p=1$. The child $x_1=0$ learns $[x_2=0]+[x_3=1]+[x_4=1]\le0$,
with infinite weight, so $\Phi=4$: the variables $x_2=1$, $x_3=0$, $x_4=0$
are forced and $x_5,x_6$ are free. The candidates in this child:
\begin{itemize}
 \item $x_2$ has gap $+2$ over two free neighbours, so $p=0$. Its other child
 is cut at once, and the measure splits as $4+0$: $I=0$, because the
 constraint had already forced $x_2$.
 \item $x_3$ has gap $0$, so its preferred value is $1$, which the learned
 constraint forbids. The other child learns $[x_2=1]+[x_6=1]\le0$, which
 contradicts $x_2=1$. Both children have soft count $0$: the whole node is
 cut, because $x_3=0$ is dominated whenever $x_2=1$.
 \item $x_4$ has gap $0$; its preferred child is forbidden, and the other
 child learns $[x_5=1]+[x_6=1]\le0$. The split is $0+1$, so $I=2$ bits.
 \item $x_5$ and $x_6$ each give the split $2+1$, so $I=\log_2(4/3)=0.415$
 bits.
\end{itemize}
The adaptive rule takes $x_3$ and cuts the node. In a fixed order starting
with $x_2$, this step gains nothing.
\end{example}

\section{Small searches}

The search branches by the adaptive rule or in index order. A new constraint
gets its best weight, vacuous constraints get weight $0$, and a node is cut
when a constraint is violated or $\Phi<1$. Every search returned the optimum.

\begin{center}
\begin{tabular}{rrrrr}
\toprule
$n$ & $d$ & no learning & fixed order & adaptive\\
\midrule
10 & 3 & 2047 & 57 & 21\\
10 & 5 & 2047 & 81 & 37\\
12 & 4 & 8191 & 65 & 25\\
12 & 6 & 8191 & 99 & 29\\
12 & 8 & 8191 & 241 & 81\\
14 & 6 & 32767 & 219 & 101\\
16 & 6 & 131071 & 623 & 89\\
16 & 10 & 131071 & 991 & 255\\
\bottomrule
\end{tabular}
\end{center}
Here $d$ is the expected degree of a random formula with one clause per pair
and random signs, and the entries are numbers of search-tree nodes. The
aligned circulant formula $C_{9,2}$ needs $429$ nodes in both orders, with
$42$ leaves against $70$ for domination alone. The information collected along
a root-to-leaf path is only $0.5$ to $5.2$ of its $9$ bits: on aligned
formulas new constraints are largely redundant, and the measure shows it.

\section{What comes next}

The running time is bounded by the total soft count at the leaves. By
additivity, this total is $2^n$ times the expected value of $2^{-\sum_tI_t}$
along one random root-to-leaf walk. The walk moves to a child with
probability proportional to the child's soft count. With weights fixed at
creation, this walk is explicit: at each step the branching variable takes
value $b$ with probability $\pi_x(b)$, and the new constraint's factor follows
Lemma~\ref{lem:factor}. So a bound $c^n$ with $c<2$ for a class of formulas
reduces to the following. Along this walk, show that the information
$\sum_tI_t$ is at least a constant times $n$, except with exponentially small
probability. That is the next step.

\paragraph{Checks.} \texttt{explore\_soft\_count.py}, each mode in seconds:
\begin{itemize}
 \item \texttt{examples} gives Examples~\ref{ex:A}--\ref{ex:C}.
 \item \texttt{search} gives the table and checks, at every branching step,
 additivity and Lemma~\ref{lem:factor} ($3{,}387$ steps). At sampled nodes it
 compares the closed form with the sum over completions and checks the bound
 of Proposition~\ref{prop:basic} ($358$ nodes). It also checks the optimum.
 \item \texttt{reopt} re-optimizes all weights from scratch at $170$ sampled
 steps and confirms
 $\Phi^*(x=0)+\Phi^*(x=1)\le\Phi^*(\text{parent})$. It also measures that
 re-optimization lowers $\ln\Phi$ by $2.7$ on average.
 \item \texttt{aligned} runs $C_{9,2}$.
\end{itemize}

\end{document}
